\scanpage{280}
$\xi$ de $\uline{\deleted{\unsure{\operatorname{Pol}}}\inserted{\operatorname{N\acute{e}r}}}'(X_T/T)$ et $\xi'$ de $\uline{\deleted{\unsure{\operatorname{Pol}}}\inserted{\operatorname{N\acute{e}r}}}'(X_{T'}/T')$ définissent la $\unsure{=}$ prépolarisation sss leurs images inverses dans $\uline{\deleted{\unsure{\operatorname{Pol}}}\inserted{\operatorname{N\acute{e}r}}}'(X_{T\times_S T'}/T\times_S T')$ sont égales, i.e. si \unsure{l'i} $p_1(\uline{L})p_2(\uline{L}')^{-1}$ est \deleted{\unsure{torsion.}} équivalent à $0$ rel. à $T\times_S T'$.

Si p. ex. $\uuline{\operatorname{Pic}}^{\tau}_{X/S}$ est plat sur $S$ et si le quotient $\inserted{\uuline{\deleted{\unsure{\operatorname{Pol}}}\inserted{\operatorname{N\acute{e}r}}}_{X/S}}\ \uuline{\operatorname{Pic}}_{X/S}/\uuline{\operatorname{Pic}}^{\tau}_{X/S}$ existe ``raisonnablement'' [i.e. le \deleted{\unsure{grp}} \unsure{graphe} d'équivalence défini par $\uuline{\operatorname{Pic}}^{\tau}_{X/S}$ est effectif, et la projection $\uuline{\operatorname{Pic}}_{X/S}\to\uuline{\operatorname{Pic}}_{X/S}/\uuline{\operatorname{Pic}}^{\tau}_{X/S}$ fid plate] -- p. ex. si $\uuline{\operatorname{Pic}}^{\tau}_{X/S}$ est \uline{propre} sur $S$ -- alors \deleted{\unsure{on prouve}} $\uline{\deleted{\unsure{\operatorname{Pol}}}\inserted{\operatorname{N\acute{e}r}}}''_{X/S}$ s'identifie au foncteur associé à \deleted{\unsure{Alors}} $\uuline{\operatorname{N\acute{e}r}}_{X/S}$, donc les néronisations de $X/S$ sont les sections de $\uuline{\operatorname{N\acute{e}r}}_{X/S}$ sur $S$, ce qui justifie la terminologie et \deleted{\unsure{en outre}} la notation

\deleted{Deux}
\deleted{Deux néronisations sont dites \uline{équivalentes} si (on appelle néronisation)}

On définit un ss-ensemble $\operatorname{N\acute{e}r}'{}^+(X/S)$ de $\operatorname{N\acute{e}r}'(X/S)$ par les classes de faisceaux inv. sur $X$ \deleted{\unsure{abs}} amples rel. à $S$, d'où un ss-foncteur $\uline{\operatorname{N\acute{e}r}}^{\unsure{+}\prime}_{X/S}$, donnant naissance à un ss-foncteur $\uline{\operatorname{N\acute{e}r}}^{+\prime\prime}_{X/S}$ de $\uline{\operatorname{N\acute{e}r}}''_{X/S}$, qui s'explicite en remplaçant ci-dessus les ``faisceaux inversibles'' par ``faisceaux inversibles rel! amples''. Dans le cas particulier
\mspage{279}
